\chapter{Conclusions and Future Work}
\label{conclusions}

\section{Conclusions}

This thesis addressed the practical problem of bridging declarative SVG input and explicit browser GPU rendering workflows.
The resulting \texttt{Svg2GPU} system delivers a coherent conversion pipeline built around typed parsing, style normalization, geometry conversion, and render integration.

The central outcome is that a robust SVG-to-GPU baseline can be implemented within bachelor-thesis constraints when scope is explicit and architecture is modular.
The project does not claim full SVG feature parity; instead, it prioritizes deterministic behavior and reproducible validation for a meaningful supported subset.

\section{Main achieved results}

The work produced the following concrete results:

\begin{itemize}
\item a parser-oriented typed model for supported SVG primitives and styles;
\item a path-focused conversion flow suitable for fill/stroke rendering pipelines;
\item practical integration with browser-oriented rendering classes;
\item interactive verification support through playground and documentation tooling;
\item executable validation for core functionalities F1 and F2.
\end{itemize}

These results demonstrate both technical feasibility and engineering discipline.

\section{Limitations}

The current version remains intentionally bounded:

\begin{itemize}
\item partial SVG specification coverage,
\item limited depth for advanced effects and full text handling,
\item triangulation improvements still possible for difficult contour classes,
\item WebGL-first practical execution with future WebGPU alignment path.
\end{itemize}

These limitations are aligned with project scope and do not invalidate achieved objectives.

\section{Future work}

The most relevant continuation directions are:

\begin{enumerate}
\item extend command and primitive coverage while preserving parser determinism,
\item strengthen contour decomposition and pathological-case triangulation handling,
\item consolidate backend abstraction for cleaner WebGPU-native execution,
\item add benchmark-driven optimization across realistic datasets,
\item expand automated regression suites for long-term maintainability.
\end{enumerate}

\section{Roadmap prioritization}

Not all continuation directions have equal impact.
For the next iteration, the recommended order is:

\begin{enumerate}
\item parser and conversion hardening for edge-case robustness,
\item triangulation and stroke-quality improvements,
\item performance profiling and optimization on representative datasets,
\item backend evolution toward cleaner WebGPU-native operation.
\end{enumerate}

This prioritization keeps the project risk-aware.
It strengthens correctness first, then improves output quality and performance, and finally addresses backend modernization.

\section{Final remark}

The practical value of this thesis is the conversion boundary itself: turning SVG semantic data into explicit, testable, and renderable GPU-oriented structures.
That boundary is the key enabler for future performance and feature evolution, and \texttt{Svg2GPU} provides a solid foundation for that path.
